% Keep the temporary build directory out of the PDF trailer ID.
\ifdefined\pdfvariable\pdfvariable suppressoptionalinfo 512\relax\fi
\documentclass{ltjbook}
\title{紙塔の組立帳}
\author{段丘模型工房}
\date{}
\begin{document}
\frontmatter
\maketitle
\tableofcontents
\mainmatter
\chapter{基壇}
\section{台の角}
図\ref{fig:base}の型に沿って灰色の台紙を折る。
\begin{figure}[htbp]
\centering\rule{34mm}{4mm}
\caption{基壇の型}\label{fig:base}
\end{figure}
\newpage
\section{台紙の裏}
返した紙の端に塔の向きを書き、折り筋を表まで通す。
\chapter{上層}
\section{高さの比較}
表\ref{tab:height}の二つの層を順に積み、塔の影が台の中央にくるか調べる。
\begin{table}[htbp]
\centering\caption{各層の高さ}\label{tab:height}
\begin{tabular}{lr}層 & 高さ（ミリ）\\下段 & 12\\上段 & 8\end{tabular}
\end{table}
\appendix
\chapter{保管箱}
外した部材は台札\cite{box}の順に箱へ戻す。
\backmatter
\begin{thebibliography}{9}
\bibitem{box} 段丘模型工房『紙塔の台札』。
\end{thebibliography}
\end{document}
